\section{Bonds: features, legal terms, markets and issuers}

These entries describe what a bond promises, the contract that holds the promise, the legal and tax rules that shape issuance and trading, and how governments, agencies, supranational lenders and companies borrow. Pricing is in \hyperref[meth:bond-price-from-yield]{pricing from a yield}; trading protocols and dealer intermediation are in \hyperref[def:fixed-income-market-structure]{fixed-income market structure}. Legal and tax rules are U.S. federal law as of 18~October 2026 unless stated.

\begin{qfdefinition}[Bond features and the indenture: issuer, maturity, par, coupon, currency, seniority and security; the trustee; affirmative, negative, maintenance and incurrence covenants, and investment-grade against high-yield covenant packages]\label{def:bond-features-indenture}
\qffield{Statement} A bond obliges the issuer to pay interest on stated dates and repay principal at maturity; failure to pay, or to meet other terms, is legal default \citep[chap.~10, p.~235]{s7-fabozzifixed}. Its features are the issuer, the maturity date (term to maturity is the time remaining), the par or face value on which the coupon is computed and which is repaid, the coupon rate and its frequency, the currency of payment, and the bond's rank and security. Maturity matters because it fixes the life of the cash flows, the point on the yield curve that prices the bond and much of its price volatility; the larger the coupon, the less the price moves for a given change in rates (chap.~1, pp.~4--6).

\emph{Indenture and trustee.} ``The legal document that sets forth all of the parameters of the agreement between the bondholders and the issuer is the indenture'' (chap.~1, p.~4). It sets three things: payments to bondholders, limits on issuer behaviour that could harm repayment, and the enforcement mechanisms if the issuer does not perform (chap.~40, p.~994). Because one holder cannot easily monitor or enforce it, a corporate trustee acts in a fiduciary capacity for holders; the trustee is paid by the issuer and does only what the indenture provides, often relying on the issuer's certificates (chap.~10, pp.~236--237).

\emph{Covenants.} Loan agreements are written as representations and warranties, \emph{affirmative} covenants (things the borrower must do) and \emph{negative} covenants (things it must not do). They are tested either continuously, \emph{maintenance} tests of leverage or coverage, or only when the borrower takes an action such as a buyback or special dividend, \emph{incurrence} tests (chap.~11, p.~266). Bond covenants are mostly negative and incurrence-based. A Debt covenant permits new borrowing only if a pro forma coverage or leverage ratio test is met; Liens and negative-pledge covenants keep other debt from being secured ahead of the bonds; Restricted Payments, Asset Sales and Mergers covenants limit dividends, disposals and asset transfers; Change of Control gives holders a put at 101 (chap.~40, pp.~997--1000).

\qffield{Conditions} \emph{Investment grade against high yield.} Investment-grade bonds usually have no subsidiary guarantees, rarely a Debt, Restricted Payments or Asset Sales covenant, a limited negative pledge with a secured-debt basket of 10 to 15 percent of a balance-sheet measure, and since 2010 usually a Change of Control put that also requires a downgrade below investment grade. High-yield bonds typically carry guarantees, Debt, Restricted Payments, Asset Sales and Change of Control covenants with many carve-outs (chap.~40, pp.~997--1000). Covenants vary over time with investors' concern for event risk, and by 2020 more than 80 percent of large U.S. and European corporate loans had no maintenance covenant (pp.~995--997). The prospectus describes the covenants, but the indenture controls (p.~997).

\qffield{Practice} Interpretation: read a bond's features as answers to four questions. Who pays (issuer, guarantors, and whether the bond ranks behind a subsidiary's own creditors)? When and how much (coupon, frequency, currency, maturity)? Under what protection (security, covenants, trustee powers, the 25 percent of holders usually needed to accelerate, p.~996)? And what can change the cash flows (\hyperref[def:bond-cash-flow-contingencies]{contingency provisions})?

\qffield{Pitfalls} Treating ``secured'' or ``senior'' as a guarantee of recovery: since the 2000s high-yield issuers have layered first-, second- and intermediate-lien debt, so recovery ``depends on the covenants'' (pp.~1001--1002). Assuming that a loan's maintenance covenants protect bondholders, when lenders can waive them (p.~996). Confusing the trustee's duties with active monitoring.

\qffield{See also} \hyperref[def:bond-cash-flow-contingencies]{cash-flow structures and contingencies}; \hyperref[meth:corporate-funding-ig-hy]{investment-grade and high-yield funding}; \hyperref[mech:cds-contract-credit-events]{credit default swaps}.

\end{qfdefinition}

\begin{qfdefinition}[Bond cash-flow structures and contingency provisions: bullet, sinking-fund, serial and amortising repayment; fixed, zero, step-up, deferred, payment-in-kind, floating and inflation-linked coupons; call, make-whole, put and conversion rights, and who benefits from each]\label{def:bond-cash-flow-contingencies}
\qffield{Statement} \emph{Principal.} Most corporate bonds are term (bullet) bonds repaying all principal at maturity. A sinking fund requires the issuer to retire part of the issue each year, by buying bonds in the market when they trade below par or by calling bonds chosen by lot at the sinking-fund price, usually par. Serial bonds and equipment trust certificates mature in instalments, each priced off its own point on the curve \citep[chap.~1, pp.~5 and~13; chap.~10, p.~244]{s7-fabozzifixed}. A fully amortising loan or mortgage pays principal in every instalment.

\emph{Coupons.} Fixed coupons are set near the market yield at issue so the bond prices near par; original-issue discount bonds and zero-coupon bonds pay less or no coupon and accrete toward par (chap.~1, pp.~6--7; chap.~10, pp.~238--239). Step-up notes raise the coupon on a schedule; deferred-interest bonds pay nothing for an initial period; payment-in-kind bonds may pay interest in more bonds (chap.~10, pp.~240 and~257). A floater pays a reference rate plus a quoted margin, possibly with a cap or floor; an inverse floater moves against the reference rate. Inflation-linked bonds index principal, coupons or both to a price index, with a lag (chap.~1, pp.~7--9).

\emph{Contingency provisions.} A call lets the issuer redeem early at a scheduled price, usually after a protection period; a non-refundable bond may still be redeemed with funds that do not come from cheaper new debt, so it protects less than a non-callable one. A make-whole call pays the greater of par and the present value of the remaining payments at a Treasury yield plus a small spread, so its price rises as rates fall. A put lets the holder sell back at par on stated dates. A convertible bond can be exchanged for a fixed number of the issuer's shares, an exchangeable bond for another company's shares (chap.~1, pp.~11--14; chap.~10, pp.~246--247).

\qffield{Conditions} \emph{Who benefits.} The call and the sinking fund's optional acceleration benefit the issuer, which exercises them when rates have fallen and the bond trades above the call or sinking-fund price; the holder then reinvests at lower yields and loses price upside, so callable bonds yield more (chap.~1, pp.~11--14). The put and conversion benefit the holder, who pays for them with a lower yield. A make-whole call transfers little value from holders when rates fall, because its price rises with the bond's value (p.~12). High-yield indentures add an equity clawback: up to 35 percent of the issue may be redeemed during the first three years, at par plus one year's coupon, from the proceeds of a share offering (chap.~10, p.~258).

\qffield{Practice} Interpretation: classify a structure by what can change the timing or size of the cash flows and who controls the change. Fixed-schedule bonds can be priced from a curve (\hyperref[meth:bond-price-from-yield]{pricing}). Bonds whose cash flows depend on rates through an option need a model and effective rather than yield duration (\hyperref[meth:curve-based-empirical-duration]{curve-based measures}).

\qffield{Pitfalls} Quoting yield to maturity for a premium callable bond; the yield to the relevant call is lower (\hyperref[def:bond-yield-spread-measures]{yield to worst}). Treating a step-up or PIK bond's initial coupon as its cost of debt. Ignoring that sinking-fund purchases at par can retire a bond trading at a premium.

\qffield{See also} \hyperref[def:bond-features-indenture]{features and indenture}; \hyperref[def:tips-breakevens]{TIPS}; \hyperref[meth:floating-money-market-yields]{floating-rate notes}.

\end{qfdefinition}

\begin{qfrule}[How law and tax shape U.S. bond issuance and trading: registration and its exemptions, shelf registration, Rule~144A, the Trust Indenture Act, municipal and federal tax treatment, original-issue discount and the settlement cycle, with an EU contrast]\label{rule:bond-issuance-legal-tax}
\qffield{Statement} \emph{Registration and its exemptions.} The Securities Act exempts from registration securities issued or guaranteed by the United States, a state or a political subdivision, and certain instrumentalities \citep[sec.~3(a)(2)]{usc-15-77c}, and notes arising from current transactions with a maturity at issue of not more than nine months \citep[sec.~3(a)(3)]{usc-15-77c}. Commercial paper relies on that exemption or on the private-placement exemption, which lifts the nine-month limit but restricts resale \citep[chap.~13, p.~302]{s7-fabozzifixed}. Eligible companies may register securities ``to be offered and sold on an immediate, continuous or delayed basis'' on Forms S-3 or F-3 (shelf registration) \citep[para.~(a)(1)(x)]{cfr-17-230-415}, which is how medium-term notes are offered continuously through agents (chap.~10, pp.~259--260). Rule~144A permits resale of unregistered securities to qualified institutional buyers, institutions owning and investing at least \$100 million in securities of unaffiliated issuers \citep[para.~(a)(1)]{cfr-17-230-144a}.

\emph{Indentures.} The Trust Indenture Act requires a qualified indenture with an eligible trustee for registered debt, but exempts securities issued otherwise than under an indenture up to an amount the Commission sets, \$50 million in any twelve months since 2015, and issues under an indenture limited to \$10 million outstanding \citep[sec.~304(a)(8)--(9)]{usc-15-77ddd} \citep[sec.~260.4a-1]{cfr-17-260-4a-1}. The Handbook's figure of \$5 million (chap.~10, p.~236) is the older threshold.

\emph{Tax.} Gross income excludes interest on state and local bonds, except private activity bonds that are not qualified, arbitrage bonds and bonds not in registered form \citep[secs.~(a)--(b)]{usc-26-103}. States may not tax U.S. government obligations or their interest, apart from non-discriminatory franchise taxes and estate or inheritance taxes \citep[sec.~(a)]{usc-31-3124}. So a municipal yield is compared with a taxable one as an equivalent taxable yield, the tax-exempt yield divided by one minus the marginal rate: 6 percent at a 40 percent rate is 10 percent; and the exemption covers coupon interest only, so a bond bought at a market discount below the \emph{de minimis} cut-off, 0.25 point a year to maturity, can produce ordinary income on its accretion \citep[chap.~9, pp.~221--223]{s7-fabozzifixed}.

\qffield{Conditions} \emph{Settlement.} A broker-dealer may not contract for settlement later than the first business day after the trade unless the parties agree otherwise; the rule excludes government and municipal securities and commercial paper \citep[para.~(a)]{cfr-17-240-15c6-1}. The T+1 cycle took effect for compliance on 28~May 2024 \citep[p.~1]{finra-regulatory-notice-23-15}. Corporate bonds therefore settle T+1 by rule, and Treasuries next day by market practice \citep[chap.~7, p.~180]{s7-fabozzifixed}.

\emph{EU contrast.} The Prospectus Regulation exempts from the prospectus obligation offers made solely to qualified investors, to fewer than 150 persons per Member State, of securities with a denomination of at least EUR~100{,}000, or requiring at least EUR~100{,}000 per investor \citep[art.~1(4)(a)--(d), consolidated text of 4~December~2024]{eu-2017-1129-prospectus}. Large denominations therefore mark the wholesale market.

\qffield{Practice} Interpretation: the rules segment the market. Registered public bonds reach all investors, need full disclosure and a qualified indenture, and trade more widely; 144A bonds reach only large institutions but come to market faster with offering memoranda, and high-yield issuers use them heavily; shelf programmes let frequent investment-grade issuers sell at short notice; tax exemption creates a separate buyer base for municipal bonds and makes in-state bonds yield less in high-tax states (chap.~9, pp.~226--227).

\qffield{Pitfalls} Comparing a municipal and a corporate yield without tax adjustment, or applying the exclusion to capital gains and market-discount accretion. Using dated thresholds; check the current rule. Treating the EU and U.S. wholesale exemptions as equivalent: one turns on denomination or investor type, the other on resale to qualified institutional buyers.

\qffield{See also} \hyperref[def:bond-market-segments-trading]{market segments}; \hyperref[def:government-agency-supranational-funding]{government funding}; \hyperref[rule:us-eu-bond-transparency]{bond transparency rules}.

\end{qfrule}

\begin{qfdefinition}[Fixed-income market segments and participants, and primary and secondary bond markets compared with equity markets]\label{def:bond-market-segments-trading}
\qffield{Statement} \emph{Segments.} The U.S. market is segmented by issuer: at the end of 2019, Treasuries \$16.7 trillion, mortgage-related \$10.3 trillion, corporate \$9.6 trillion, municipal \$3.9 trillion, federal agency \$1.8 trillion, asset-backed \$1.8 trillion and money markets \$1.0 trillion \citep[chap.~1, Exhibit~1-1, p.~4]{s7-fabozzifixed}. Bonds are further divided by maturity (money market under a year; Handbook usage short, intermediate and long), credit quality (investment grade or high yield), currency and domicile (domestic, foreign, Eurobond) and coupon type (chap.~1, pp.~5--6). \hyperref[def:fixed-income-universe]{The fixed-income universe} lists the macro trader's segments.

\emph{Participants.} Issuers are governments, agencies and government-sponsored enterprises, supranational institutions, municipalities, financial and non-financial companies and securitisation vehicles. Investors include central banks, banks (for liquidity, interest-rate management and collateral), pension funds and insurers (to fund long liabilities), mutual and money-market funds, hedge funds and individuals (chap.~16, pp.~343--344). Intermediaries are underwriters, primary dealers, interdealer brokers and electronic platforms.

\emph{Primary market.} New bonds are sold by auction (Treasuries, many sovereigns and agency discount notes), by syndicate or underwriting at a negotiated price (most corporate, many agency and some sovereign issues), by competitive bidding or negotiation for municipalities, by continuous offering through agents for medium-term notes and commercial paper, by tap or reopening of existing issues, or privately (chap.~7, pp.~173--177; chap.~8, pp.~190--192; chap.~9, p.~227; chap.~10, pp.~259--260; chap.~16, pp.~346--348).

\qffield{Conditions} \emph{Against equities.} An equity issuer usually has one class of common stock listed on an exchange and sells new shares rarely; a bond issuer may have dozens of separate, non-fungible bonds outstanding and issues often. Secondary bond trading is mainly over the counter through dealers, with quotes that are often indicative, while equities trade continuously on exchanges with firm limit-order books \citep[sec.~I, pp.~1--5]{wide-bsv-fi-survey}. Liquidity is concentrated: in the Treasury market, on-the-run issues accounted for 69 percent of interdealer volume against 27 percent for the off-the-run majority, and the most active issue, the on-the-run five-year note, averaged about \$115 billion a day in the data the chapter cites \citep[chap.~7, p.~180]{s7-fabozzifixed}; individual corporate bonds traded on fewer than half of days in one sample \citep[p.~5]{wide-bsv-fi-survey}.

\qffield{Practice} Interpretation: segment, issuance method and trading venue go together. A large, frequent, standardised issuer (the Treasury) auctions on a fixed calendar, and its bonds trade electronically at narrow spreads; a single company issues by syndicate, and its many bonds trade by request for quote or voice, with dealers holding inventory. Index construction must cope with this: bonds enter and leave and many have no daily trade (\hyperref[def:equity-fixed-income-index-types]{fixed-income indexes}).

\qffield{Pitfalls} Inferring the liquidity of a corporate bond from its issuer's equity; the issuer's shares may trade continuously while the bond does not trade for days. Using 2019 sizes and volumes as current; they are dated to the Handbook's sources.

\qffield{See also} \hyperref[def:fixed-income-market-structure]{fixed-income market structure}; \hyperref[rule:bond-issuance-legal-tax]{legal and tax rules}; \hyperref[def:government-agency-supranational-funding]{government funding}.

\end{qfdefinition}

\begin{qfmethod}[Long-term funding of investment-grade and high-yield companies: instruments, security, covenants, call terms, investors and market access compared]\label{meth:corporate-funding-ig-hy}
\qffield{Statement} \emph{Investment-grade issuers} (rated BBB$-$/Baa3 or higher) fund mainly with senior unsecured debentures, issued often and in size, many through a shelf registration or medium-term note programme \citep[chap.~10, pp.~252 and~259--260]{s7-fabozzifixed}. Their bonds carry light covenants: usually no Debt, Restricted Payments or Asset Sales covenant, a limited negative pledge, and a make-whole call rather than a fixed call schedule (chap.~40, pp.~997--1000; chap.~10, p.~247). Commercial paper and bank lines fund the short term (chap.~13, pp.~301--308).

\emph{High-yield issuers} (below BBB$-$/Baa3) include young or story companies, fallen angels and companies leveraged by restructurings or buyouts (chap.~10, pp.~256--257). They fund with secured and unsecured bonds that carry subsidiary guarantees and a full set of incurrence covenants; are typically non-callable for an initial period, then callable at declining premiums, with an equity clawback (p.~258); and may defer cash interest through deferred-interest, step-up or payment-in-kind structures (p.~257). Leveraged loans, usually secured and floating rate, sit beside the bonds (chap.~11, pp.~264--266).

\qffield{Conditions} \emph{Why the contracts differ.} Investors in high-yield debt bear more default risk and so demand tighter contracts and higher yields. Yet investment-grade holders, who lose most when an issuer is cut to high yield by a buyout or a large dividend, have the fewest covenant protections, a paradox that has lasted decades (chap.~40, p.~997). Security has become common in high yield, with 20 to 30 percent of new issues secured in a typical quarter since 2010 (p.~1001).

\qffield{Practice} Interpretation, a comparison for a funding decision:
\begin{itemize}
\item \emph{Cost and access}: investment-grade issuers borrow at small spreads, can issue opportunistically from a shelf and can roll commercial paper; high-yield issuers pay higher spreads, depend more on loan and 144A markets, and can lose access when spreads widen, as in 2008.
\item \emph{Flexibility}: investment-grade debt leaves management free to pay dividends, sell assets or add debt; high-yield debt limits these, and refinancing early costs a call premium, while a make-whole call costs the present value of the remaining payments.
\item \emph{Maturity and structure}: investment-grade issuers can sell 30-year and longer bonds; high-yield maturities are shorter (often 5 to 10 years), with deferred-coupon and secured structures for weaker credits.
\end{itemize}
Recommendation rule: a company near the boundary weighs the lower spread of an investment-grade rating against the financial policy it requires; a buyout deliberately accepts high-yield terms to fund leverage.

\qffield{Pitfalls} Reading covenant-light investment-grade bonds as low event risk; a takeover can re-lever the issuer and turn the bond into high yield (the RJR Nabisco spread widened from about 100 to 350 basis points on the 1988 buyout bid, chap.~10, p.~256). Treating ratings boundaries as fixed; they are agency judgements and change.

\qffield{See also} \hyperref[def:bond-features-indenture]{covenants}; \hyperref[rule:bond-issuance-legal-tax]{144A and shelf registration}; \hyperref[def:government-agency-supranational-funding]{government funding}.

\end{qfmethod}

\begin{qfdefinition}[Government, agency and supranational funding: sovereign auctions and instruments, domestic and external debt, municipal general-obligation and revenue bonds, agency and GSE programmes, supranational issuers, and the contrast with corporate issuance and trading]\label{def:government-agency-supranational-funding}
\qffield{Statement} \emph{Sovereigns.} The U.S. Treasury sells bills (up to one year, at a discount), notes (two to ten years), bonds (20 and 30 years), TIPS and two-year floating-rate notes by sealed-bid, single-price auction on a regular, pre-announced calendar; all accepted bids are filled at the highest accepted (stop-out) yield, no bidder may be awarded more than 35 percent of an issue, and noncompetitive bids up to \$5 million are filled in full \citep[chap.~7, pp.~172--176]{s7-fabozzifixed}. Primary dealers are expected to bid in every auction and make markets (p.~174). Other sovereigns use single- or multiple-price auctions in developed markets, and syndication, underwriting or tap sales where demand is uncertain or the market is thin (chap.~16, pp.~346--348). \emph{Domestic} debt is issued under national law in the national currency; \emph{external} debt is issued under foreign law, usually in foreign currency, and needs foreign exchange to service (p.~348).

\emph{Non-sovereign and quasi-government issuers.} Municipalities issue general-obligation bonds backed by taxing power and revenue bonds backed by a project's revenues (chap.~1, p.~4), by competitive bidding (required by most states for general obligations) or negotiation (chap.~9, p.~227). U.S. agencies and government-sponsored enterprises (Federal Home Loan Banks, Fannie Mae, Freddie Mac, Farm Credit) issue discount notes by auction or reverse inquiry, benchmark bonds by syndicate on a coordinated calendar, and callable notes to hedge mortgage prepayment; most have no explicit federal guarantee (chap.~8, pp.~185--192). Supranational lenders borrow in global capital markets: the IBRD has issued since 1947, IDA since 2018, both rated Aaa/AAA, through benchmark, non-core-currency, structured and discount notes, with 61 percent of IBRD funding in U.S. dollars in fiscal 2026 \citep[pp.~1--2]{world-bank-ibrd-ida-factsheet-2026}.

\qffield{Conditions} \emph{Against corporate issuance and trading.} Governments issue regularly, in large benchmark sizes and by transparent auction, to keep a liquid curve; companies issue when they need funds, mainly by underwritten syndicate, in many smaller, non-fungible bonds. Treasuries trade in a deep dealer and electronic interdealer market at bid-ask spreads of a fraction of a thirty-second, and serve as the pricing and hedging benchmark; corporate bonds trade less and at wider spreads (chap.~7, pp.~178--182) \citep[pp.~1--5]{wide-bsv-fi-survey}. Agency debt sits between: primary dealer volume averaged 9.5 percent of Treasury volume in 2010--2020 (chap.~8, p.~192).

\qffield{Practice} Interpretation: credit and currency define the risk. A sovereign borrowing in its own currency can always pay nominally but may inflate; in foreign currency it can default for lack of reserves (chap.~16, pp.~348--349; \hyperref[def:sovereign-credit-defaults]{sovereign credit}). Municipal credit depends on the pledged revenue, and tax exemption shapes its buyers (\hyperref[rule:bond-issuance-legal-tax]{tax}). Agencies trade at a small spread to Treasuries for credit and liquidity; Fannie Mae and Freddie Mac have been in conservatorship since September~2008 without an explicit guarantee of their debentures (chap.~8, p.~187).

\qffield{Pitfalls} Treating agency or supranational debt as Treasury debt; their spreads widened in 2008. Using 2019 auction sizes and schedules as current; Treasury changes them (TIPS calendar in \hyperref[def:tips-breakevens]{TIPS}).

\qffield{See also} \hyperref[def:bond-market-segments-trading]{market segments}; \hyperref[meth:corporate-funding-ig-hy]{corporate funding}; \hyperref[meth:floating-money-market-yields]{bill quotation}.

\end{qfdefinition}
